%%%%%%%%%%%%%%%%%%%%%%%%%%%%%%%%%%%%%%%%%
% Seismica Submission Template
% Falcon 9 fireball paper working manuscript -- version 12
%%%%%%%%%%%%%%%%%%%%%%%%%%%%%%%%%%%%%%%%%

\documentclass[]{seismica}
\usepackage{tabularx}

% All manuscript graphics are stored in the figures subdirectory.
\graphicspath{{figures/}}

\title{Near-Field Seismo-Acoustic Observations of the 2016 Falcon 9 Explosion at Kennedy Space Center}
\shorttitle{Near-field observations of the 2016 Falcon 9 explosion}

% TODO: add ORCID identifiers and confirm corresponding-author address.
\author[1]{Glenn Thompson\thanks{Corresponding author: EMAIL TO BE CONFIRMED}}
\author[2]{Robert G. Brown}
\author[1]{Alexandra Farrell}
\author[2]{John Kiriazes}
\author[1]{Stephen R. McNutt}
\author[1]{Christopher Mehta}
\affil[1]{School of Geosciences, University of South Florida, Tampa, Florida, USA}
\affil[2]{NASA Kennedy Space Center, Florida, USA (current affiliations to be confirmed)}

% TODO: complete and verify CRediT contributions.
\credit{Conceptualization}{Glenn Thompson, Stephen R. McNutt}
\credit{Methodology}{Glenn Thompson}
\credit{Software}{Glenn Thompson}
\credit{Formal Analysis}{Glenn Thompson}
\credit{Investigation}{Glenn Thompson, Robert G. Brown, Alexandra Farrell, John Kiriazes, Stephen R. McNutt, Christopher Mehta}
\credit{Visualization}{Glenn Thompson}
\credit{Writing - Original draft}{Glenn Thompson}
\credit{Writing - Review \& Editing}{All authors}

\begin{document}

\makeseistitle{
  \begin{summary}{Abstract}
  On 1 September 2016, a Falcon~9 launch vehicle was destroyed during
  preparations for a pre-launch static-fire test at Space Launch Complex~40
  (SLC-40), Cape Canaveral Air Force Station (CCAFS; now Cape Canaveral Space
  Force Station), Florida. The accident was recorded approximately
  \SI{1.4}{\kilo\metre} away by a colocated broadband seismometer and
  three-element infrasound array operated by the University of South Florida
  near the Astronaut Beach House on adjacent Kennedy Space Center (KSC)
  property, providing an unusually close calibrated seismo-acoustic record of
  a major technological explosion.
  %
  The observations were analyzed as an independent forensic record of the
  accident: to test for detectable energetic precursors, reconstruct the
  chronology, quantify event amplitudes, determine whether signals originated
  consistently from the launch complex, and evaluate their propagation
  characteristics. Manual reconciliation identified 153 impulsive infrasound
  events distributed among four phases over approximately \SI{26}{\minute}.
  A simple positive-compression/negative-rarefaction morphology was measurable
  for 152 events, while the principal explosion was retained as a complex
  noncanonical waveform. No
  anomalous seismic or infrasonic precursor was detected during the preceding
  \SI{37}{\hour}. The initial upper-stage event occurred at 13:07:11.913~UTC;
  approximately \SI{3.60}{\second} later, the principal explosion generated a
  median sensor-specific positive pressure of \SI{1.392}{\kilo\pascal}, a
  peak-to-peak pressure of \SI{1.549}{\kilo\pascal}, and the strongest ground
  motion of the sequence. A prolonged interval of
  continuous signal followed the initial phase, but its physical origin
  remains uncertain.
  %
  High-quality array solutions pointed toward SLC-40 and gave a
  coherence-weighted apparent velocity of \SI{351.98}{\metre\per\second}, in
  close agreement with the meteorologically predicted
  \SI{350.98}{\metre\per\second}. Across a measurement-quality subset of 46
  events, peak-to-peak pressure and three-component peak ground velocity were
  nearly proportional (log--log slope 0.962; $R^2=0.953$). Seismic records were
  dominated by ground-coupled airwaves, and classical N-wave morphology was
  not observed.
  A single near-field seismo-acoustic array therefore provided an independent
  reconstruction of a complex launch-vehicle failure without reliance on
  proprietary engineering telemetry.
  \end{summary}
  %
  \begin{summary}{Non-technical summary}
  A seismic and low-frequency sound monitoring station approximately
  \SI{1.4}{\kilo\metre} from the launch pad recorded the 2016 Falcon~9
  accident at SLC-40. The launch pad was at Cape Canaveral Air Force Station,
  while the monitoring station was located on adjacent Kennedy Space Center
  property. The recordings provide an independent account of what happened:
  whether any detectable energetic signal preceded the visible failure, how
  many subsequent events occurred, how large they were, where they originated,
  and how quickly their pressure waves travelled. No unusual precursor was
  detected during the preceding \SI{37}{\hour}. The accident was not a single
  explosion, but a sequence of 153 impulsive pressure events distributed
  among four phases over approximately \SI{26}{\minute}; 152 were represented
  by the standard short-pulse measurement algorithm, while the principal
  explosion had a longer, more complex waveform. By comparing arrival times and waveforms at
  three infrasound sensors, we show that the events consistently originated
  from the direction of the launch complex. The study demonstrates how a
  small, continuously recording geophysical station can provide an independent
  forensic record of a complex launch-vehicle failure.
  \end{summary}
}


\section{Introduction}
\label{sec:introduction}

Rocket launches generate some of the strongest routinely occurring
anthropogenic seismic and infrasonic signals on Earth. During nominal launches,
intense acoustic radiation is produced principally by the interaction of rocket
exhaust with the atmosphere and launch infrastructure, and these airwaves can
couple efficiently into the shallow ground. Rocket-generated signals have been
studied for both operational monitoring and geophysical research, including
global surveys of infrasonic signatures from routine space launches
\citep{pilger_rocket_2021}. Catastrophic launch failures remain rare,
particularly when recorded at close range by calibrated seismic and infrasound
instruments.

On 1 September 2016, a Falcon~9 launch vehicle was destroyed during
preparations for a pre-launch static-fire test at Space Launch Complex~40
(SLC-40), Cape Canaveral Air Force Station (CCAFS; now Cape Canaveral Space
Force Station), Florida. The vehicle and its AMOS-6 payload were lost, and the
launch complex sustained extensive damage. The accident was documented by
several publicly available videos and was subsequently investigated by SpaceX,
providing a high-energy technological explosion with known source location,
timing, and geometry for which the geophysical record can be compared with
independent visual observations and a subsequently released engineering
chronology.

In the immediate aftermath, the USF recordings were examined as an independent
forensic record of the accident. The principal questions were whether any
detectable energetic precursor preceded the visible failure; how many discrete
events occurred and when; how large the principal events were; whether the
signals originated consistently from SLC-40; and whether their apparent
propagation speeds were compatible with atmospheric airwaves under the observed
meteorological conditions. These questions motivated the original analysis
conducted for SpaceX and the subsequent scientific investigation.

At the time of the accident, the University of South Florida (USF) was
operating the BCHH seismo-acoustic array near the Astronaut Beach House on
Kennedy Space Center property, approximately 1.4~km north-northeast of SLC-40.
The installation comprised a colocated broadband seismometer and three-element
infrasound array and formed part of a broader USF program that used rocket
launches to evaluate seismo-acoustic instrumentation between volcanic field
deployments, train students in field operations, and investigate near-field
rocket-generated wavefields. The BCHH array fortuitously recorded the complete
accident sequence with high signal-to-noise ratio, without reliance on vehicle
telemetry or a triggered recording system.

Previous seismo-acoustic studies of large technological explosions have
generally relied on regional observations. Examples include the 2005 Buncefield
oil-depot explosion \citep{ottemoller_2008,ceranna_2009,green_2008}, the 2014
Moerdijk chemical-plant explosions \citep{smink_2019}, and the 2020 Beirut
ammonium-nitrate explosion \citep{pilger_beirut_2021,kim_pasyanos_2022}.
Atmospheric propagation modelling was central to interpreting arrival times,
amplitudes, source location, and explosive yield in those studies. By contrast,
the 1.4-km source distance places the BCHH array in the near field, allowing
direct-arrival analysis with comparatively little influence from atmospheric
stratification or ducting, while still requiring explicit treatment of local
wind and finite-range geometry. This proximity allows the chronology and
relative strength of more than 150 individual impulsive events to be resolved
and compared with independent video observations. The dataset therefore
provides an unusually detailed view of the evolution of a catastrophic
launch-vehicle failure.

The objectives of this study are to:
\begin{enumerate}
  \item test the continuous seismic and infrasound record for detectable
        energetic precursors to the visible failure;
  \item construct a quantitative chronology of the accident and determine the
        number and temporal distribution of discrete impulsive events;
  \item measure the acoustic and seismic amplitudes of the principal events
        and characterize their associated wavefields;
  \item determine whether the impulsive arrivals originated consistently from
        the direction of SLC-40 and whether their apparent propagation speeds
        were compatible with the surveyed geometry and observed atmospheric
        conditions; and
  \item evaluate the value and limitations of near-field seismo-acoustic
        observations for independent forensic analysis of launch-vehicle
        accidents and other complex technological explosions.
\end{enumerate}

This study does not attempt to determine the internal engineering mechanism
that initiated the vehicle failure. In particular, the absence of a detectable
energetic precursor cannot exclude non-energetic or seismo-acoustically
undetectable processes. Instead, the study focuses on what can be inferred
directly from the near-field geophysical observations and on how those
observations complement the publicly available visual record and the
subsequently released failure investigation.


\section{Instrumentation and data}
\label{sec:data}

\subsection{BCHH seismo-acoustic array}

The BCHH array was initially deployed on 24 February 2016 as part of a
three-station USF network on Kennedy Space Center property, but two stations
were removed in May 2016. The network was established to evaluate and
recalibrate seismic and infrasound instruments returned from the 2015
Sakurajima experiment, during which anomalous amplitude responses had been
identified. It subsequently developed into a longer-term program in which
equipment was kept in continuous operational use between volcanic field
campaigns. This supported routine instrument-performance checks, student
training in equipment deployment and maintenance, teaching, and research on
rocket-generated seismic and acoustic wavefields. The station and channel history, 
including sensor positions and response metadata, is preserved in the network 
metadata and StationXML files used during waveform processing.

The 1 September 2016 Falcon~9 accident at SLC-40, Cape Canaveral Air Force
Station (CCAFS), was recorded at BCHH near the Astronaut Beach House on
adjacent Kennedy Space Center (KSC) property
(Figure~\ref{fig:site_array}). Located approximately 1.42~km
north-northeast of SLC-40, BCHH provided favorable geometry for recording
both routine launches and, fortuitously, the accident. At the time of the
event, the array comprised a three-component Nanometrics Trillium Compact
Posthole broadband seismometer near the center of a three-element infrasound
array with an aperture of approximately 30~m. The seismometer was buried
approximately 0.5~m below ground level to improve thermal stability and reduce
wind-induced noise.

The three infraBSU differential-pressure sensors were mounted individually on
short stakes approximately 20--30~cm above ground level in an approximately
equilateral triangular configuration around the seismometer. All six channels
were digitized continuously at 250 samples~s$^{-1}$ using a GPS-synchronized
Nanometrics Centaur recorder. 

\begin{figure*}[t]
  \centering
  \includegraphics[width=0.96\textwidth]{fig01_slc40_bchh_location.pdf}
  \caption{Location and geometry of the BCHH seismo-acoustic array.
  (a) Map of the southern Kennedy Space Center--Cape Canaveral launch area
  showing the USF BCHH array on Kennedy Space Center property, approximately
  1.4~km north-northeast of Space Launch Complex~40 (SLC-40) at Cape
  Canaveral Air Force Station; Space Launch Complex~41 is shown for
  geographic reference.
  (b) Geometry of the BCHH array, comprising a three-component broadband
  seismometer near the center of a three-element triangular infrasound array
  with an aperture of approximately 30~m. Pressure-sensor positions were
  reconstructed from cable-end locations measured after the sensors had been
  removed; metre-scale relative-position uncertainty is therefore possible.}
  \label{fig:site_array}
\end{figure*}


\subsection{Surveyed geometry}

The array geometry was reconstructed on 2 November 2016 after the pressure
sensors had been removed. The surveyed points therefore represent cable ends
rather than the original sensor diaphragms. Relative sensor-position errors of
approximately 1--2~m are plausible and are more realistic than the formal
centimetre-scale precision of the survey itself. Using these coordinates and
the known launch-pad location, source-to-sensor distances were 1439.96~m for
DD1, 1408.34~m for DD2, and 1412.85~m for DD3. The source-to-seismic-site
distance was approximately 1419.9~m. The WGS84 source-to-array azimuth was
19.436$^\circ$, corresponding to an expected back azimuth of
199.436$^\circ$ at BCHH (Table~S2). These coordinates formed the basis for
the differential-travel-time corrections and planar- and finite-range array
analyses described in Section~\ref{sec:methods}.


\subsection{Instrument responses and pressure calibration}

The broadband-seismometer response was removed using the manufacturer
calibration for the Trillium Compact Posthole--Centaur instrument chain. The
infraBSU channels were converted to pressure using empirical, sensor-specific
field calibration factors derived from comparative tests across the set of 25
infraBSU sensors operated by USF. Available calibration records indicate that
these experiments included colocated lift tests and impulsive tests such as
door slams, allowing relative differences in sensor sensitivity to be
quantified. The adopted BCHH calibration factors are summarized in
Table~\ref{tab:instrument_calibration}; their implementation during response
correction is described in Section~\ref{sec:methods}.

The original analysis assigned an approximate absolute pressure uncertainty of
$\pm20\%$. Because the basis for this value is not encoded in the modern
processing notebooks, it will be verified against the original calibration
records before final submission. Timing uncertainty reflects the 4-ms sampling
interval, GPS synchronization, and manual arrival picking. For clear impulsive
arrivals, absolute pick uncertainties are estimated to be better than
approximately 0.02~s.

\begin{table}
\centering
\caption{Instrument channels, sampling rate, and calibration factors used for
response correction. The infraBSU values are empirical field sensitivities.}
\label{tab:instrument_calibration}
\begin{tabular}{lll}
\hline
Instrument/channel & Sampling rate & Sensitivity \\
\hline
Trillium Compact DHZ/DHN/DHE
  & 250~Hz
  & $3.0172\times10^8$ counts $(\mathrm{m\,s^{-1}})^{-1}$ \\
infraBSU DD1
  & 250~Hz
  & 8.10 counts~Pa$^{-1}$ \\
infraBSU DD2
  & 250~Hz
  & 0.690 counts~Pa$^{-1}$ \\
infraBSU DD3
  & 250~Hz
  & 10.0 counts~Pa$^{-1}$ \\
\hline
\end{tabular}
\end{table}


\subsection{Meteorological observations}

Meteorological observations were obtained from the KSC weather-tower network
during the accident sequence. The network comprised 68 instrumented towers
spanning measurement heights of approximately 4--150~m above ground level.
Available variables included wind speed, wind direction, air temperature,
relative humidity, and barometric pressure at 5-min intervals between 13:05
and 13:35~UTC. These observations provided unusually dense spatial and vertical
sampling of atmospheric conditions during the accident and were used to
estimate the still-air sound speed, the along-path wind correction, and the
expected acoustic travel time between SLC-40 and BCHH
(Section~\ref{sec:methods}).


\subsection{Publicly available video and audio}

The seismic and infrasound observations were compared with publicly available
video recordings of the accident. These recordings provide an independent
visual chronology and allow selected geophysical observations to be related to
visible stages including the initial flash, growth of the fireball, structural
collapse, descent of vehicle or payload components, and subsequent energetic
activity. Public availability makes the source material independently
inspectable, although several recordings lack reliable absolute timestamps and
their camera positions, frame rates, and timing uncertainties require further
documentation.

The videos were used principally to interpret and label the geophysical
sequence. Event association, waveform amplitudes, array measurements, and the
other quantitative catalogue products were derived from the BCHH seismic and
infrasound recordings. Proprietary SpaceX footage may have been viewed during
the original forensic work, but no quantitative result reported here will rely
on material that cannot be independently inspected or redistributed; any
historical influence of such footage on an interpretation will be disclosed
explicitly. The procedures and limitations of video and audio
synchronization are described in Section~\ref{sec:methods}.


\section{Methods}
\label{sec:methods}

\subsection{Analysis overview}

The analysis was organized around the principal forensic questions posed
following the accident: whether any detectable energetic precursor preceded
the visible failure; how many discrete events occurred and when; how their
acoustic and seismic amplitudes varied; whether they originated consistently
from SLC-40; and whether their apparent propagation speeds were compatible with
direct acoustic propagation under the observed atmospheric conditions.

The technical workflow comprised:
(1) correction of instrument responses and conversion to physical units;
(2) compilation of surveyed source--receiver geometry and meteorological
constraints;
(3) manual identification and reduced-time association of candidate airwave
arrivals;
(4) waveform review, pressure-pulse characterization, and generation of fixed
six-channel event segments;
(5) independent catalogue validation using continuous Bartlett beamforming;
(6) estimation of back azimuth and apparent propagation speed for all events
using time-domain planar-wave coherence analysis;
(7) assessment of finite-range effects using a horizontal circular-wavefront
model;
(8) measurement of pressure, peak ground velocity, acoustic--seismic amplitude
scaling, spectra, and pressure-to-ground transfer estimates;
(9) estimation of sound exposure, radiated acoustic energy, and illustrative
TNT-equivalent scenarios for nonoverlapping event complexes;
(10) calculation of source-oriented three-component polarization attributes
and a focused search for seismic motion preceding the ground-coupled airwave;
and (11) synchronization of public video and camera audio with the BCHH
waveforms and generation of publication and supplementary figures.

Modern waveform processing, event association, amplitude and energy
measurement, array analysis, polarization analysis, and multimedia alignment
were implemented in Python and are described directly from the executed
notebooks. These products supersede equivalent legacy MATLAB values. The
legacy MATLAB/GISMO workflow provided the original forensic display and is
retained for provenance, but no principal numerical result depends solely on
that implementation. Continuous waveform data and
metadata were originally archived in a CSS3.0 database using Antelope
\citep{brtt_antelope}, manual arrival picks were made using
\texttt{dbpick}, and the original MATLAB analysis used GISMO
\citep{thompson_gismo_2017}. The canonical Python analysis interval extended from 180~s
before to 1800~s after the opening sequence. All six waveform channels were
sampled continuously at 250~Hz. Final named-event source times, atmospheric
corrections, and figure annotations were read from shared configuration and
pipeline products rather than re-entered in individual notebooks.


\subsection{Instrument calibration and response correction}

The three seismic components were corrected using the manufacturer response
for the Nanometrics Trillium Compact Posthole seismometer and Centaur digitizer
combination. The combined sensitivity was
$3.0172\times10^8$ counts $(\mathrm{m\,s^{-1}})^{-1}$. The Centaur was
operated with a 40-V peak-to-peak input range, corresponding to a digitizer
sensitivity of $4.0\times10^5$ counts~V$^{-1}$ and an implied seismometer
sensitivity of approximately 754~V~$(\mathrm{m\,s^{-1}})^{-1}$.

The three infraBSU differential-pressure sensors were converted independently
from digital counts to pressure using empirical, sensor-specific field
calibration factors of 8.10, 0.690, and 10.0 counts~Pa$^{-1}$ for DD1, DD2,
and DD3, respectively. These factors were derived from comparative experiments
in which colocated sensors were exposed to common pressure changes and
impulsive sources. They should therefore be regarded as empirical field
calibrations rather than laboratory-traceable absolute calibrations. The
precise combination of tests used to derive the adopted factors, and the basis
for the previously assigned absolute pressure uncertainty of approximately
$\pm20\%$, will be verified from the original calibration records.

For each infrasound channel, the median count level between 20 and 2~s before
the initial explosion was removed before response correction. The empirical
pressure sensitivities were incorporated into complete instrument-response
representations while retaining the nominal poles, zeros, and digitizer
stages. Seismic traces were demeaned and linearly detrended. All traces received
a 5\% cosine taper before frequency-domain response removal.

A four-corner prefilter of 0.02, 0.05, 80, and 100~Hz and a water level of
60~dB were used to stabilize the response inversion. This prefilter was applied
only during response removal. No additional Butterworth high-pass filter was
applied during event association, amplitude measurement, or array analysis.
Corrected seismic traces were expressed as ground velocity in
m~s$^{-1}$, and corrected infrasound traces as differential pressure in Pa.


\subsection{Source--receiver geometry}

Sensor coordinates were transformed to UTM Zone~17N for array calculations.
Horizontal source--receiver distances and geographic azimuths were also
calculated independently on the WGS84 ellipsoid. The source-to-seismic-site
distance was approximately 1419.9~m, and the geometric back azimuth from the
nominal array reference toward SLC-40 was 199.436$^\circ$, measured clockwise
from north. As noted above, the pressure-sensor coordinates represent surveyed
cable ends after sensor removal, so metre-scale relative errors are plausible.

The three infrasound sensors formed an array with an aperture of approximately
30~m. Their source distances were 1439.96~m (DD1), 1408.34~m (DD2), and
1412.85~m (DD3). Individual
source-to-sensor distances were retained when calculating expected
differential acoustic travel times rather than treating the array as a single
receiver. DD2 was used as the infrasound reference channel for reduced-time
association and waveform alignment; consequently, its geometric direction to
SLC-40 differs slightly from the nominal array-center direction.

Source times for selected stages were obtained by reducing precisely reviewed
BCHH pressure arrivals using the 351~m~s$^{-1}$ source-to-array speed. Public
video was then used to interpret the corresponding visible stages; it was not
used to construct the general event catalogue. The final named source times
were 13:07:11.913~UTC (upper-stage event), 13:07:15.514~UTC (principal
explosion), 13:07:24.420~UTC (capsule impact), 13:07:24.988~UTC (capsule
explosion), and 13:07:26.458~UTC (event~024). Because the USLaunchReport video
lacks absolute timing, its time base was mapped to this geophysical chronology
rather than treated as an independent UTC reference.

\subsection{Chronology and timing conventions}
\label{sec:methods_legacy_chronology}

Three time coordinates were kept distinct. Observed UTC is the waveform time at
each receiver. Reduced arrival time corrects DD1 and DD3 only for their
differential acoustic travel times relative to DD2 and is the coordinate used
for event association. Source time subtracts the complete source-to-receiver
travel time and is used for physical chronology and multimodal figures. Thus,
the catalogue variable \texttt{reduced\_event\_epoch\_s} is an estimated DD2
arrival, not an origin time.

The meteorological prediction of 350.98~m~s$^{-1}$ was rounded to
351~m~s$^{-1}$ for source-to-BCHH travel-time reduction. The independently
measured array-scale apparent speed of approximately 352~m~s$^{-1}$ was
retained as an observational result and, where noted, used for differential
waveform alignment over the approximately 30-m aperture. The distinction is
important: the former is predicted from temperature, humidity, and along-path
wind, whereas the latter is measured from relative delays within the array.


\subsection{Meteorological correction}

Meteorological observations from the Kennedy Space Center tower network between
13:05 and 13:35~UTC were used to estimate the effective acoustic propagation
speed between SLC-40 and BCHH. The selected dataset contained 105 records from
three tower locations at seven 5-min observation times, spanning elevations of
16.5--139.9~m. Mean conditions were 26.36$^\circ$C and 83.9\% relative
humidity.

The still-air sound speed was calculated as
\begin{equation}
  c_0 = 331.3 + 0.606T + 1.26h,
  \label{eq:sound_speed}
\end{equation}
where $T$ is temperature in degrees Celsius and $h$ is relative humidity
expressed as a fraction. This yielded
$c_0=348.33~\mathrm{m\,s^{-1}}$.

Wind observations were resolved into east and north components and averaged
within 10-m height intervals. The available lower-atmosphere profile was
summarized over the nominal 0--65~m source-height interval. Because the lowest
valid wind measurements were near 40~m, this is an estimate of the lowest part
of the propagation path rather than a direct observation throughout the
surface layer. The height-averaged wind component parallel to the
SLC-40--BCHH propagation path was added to the thermodynamic sound speed:
\begin{equation}
  c_{\mathrm{eff}} =
  c_0 + \mathbf{u}\cdot\hat{\mathbf{r}}.
  \label{eq:effective_speed}
\end{equation}

The mean along-path wind component was $+2.65~\mathrm{m\,s^{-1}}$, giving
$c_{\mathrm{eff}}=350.98~\mathrm{m\,s^{-1}}$. The corresponding travel time
over 1419.9~m was 4.046~s, approximately 0.031~s shorter than the still-air
estimate. Because the propagation distance was only about 1.4~km, direct
atmospheric arrivals dominated and long-range ducting effects were not
considered. The meteorological result was rounded to 351~m~s$^{-1}$ for
source-to-BCHH reduction throughout the final workflow. It is distinct from,
but closely agrees with, the approximately 352~m~s$^{-1}$ apparent velocity
recovered independently across the infrasound array.

\subsection{Manual arrival picking and baseline correction}

Continuous waveforms were visually reviewed using Antelope \texttt{dbpick}. Candidate events were initially identified as impulsive pressure arrivals on the infrasound channels. Only non-deleted phase-$N$ picks on DD1, DD2, and DD3 were retained, and the original pick precision was preserved when exporting the arrivals to the Python workflow.

Manual review was adopted because the accident sequence occupied approximately 26~min within a modest six-channel dataset and contained closely spaced, overlapping arrivals spanning a wide range of amplitudes and waveform shapes. For this dataset, manual inspection provided a consistent reference catalogue without the iterative tuning and validation required for a new automated detector. Local seismic transients without corresponding infrasound arrivals were excluded from the airwave catalogue.

To address the forensic question of whether any detectable energetic precursor
preceded the visible failure, the complete available record from 31 August
through 1 September 2016 was also examined outside the main accident sequence.
This interval included approximately 37~h before the first accident-related
signal and 10.5~h after the final catalogued event. The search was intended to
identify anomalous impulsive seismic or infrasonic signals and cannot exclude
non-energetic, aseismic, or otherwise undetectable initiating processes.

Slow baseline variations on the response-corrected infrasound channels were removed using a centered, noncausal moving median with a 251-sample window, corresponding to approximately 1~s at 250~Hz. The estimated median baseline was subtracted from each pressure trace. Seismic traces were not subjected to this baseline operation.

The moving-median procedure was independently recomputed from the original
response-corrected stream and reproduced the canonical pressure product sample
for sample. Sensitivity tests used 0.125-, 0.25-, 0.5-, 1.0-, and 2.0-s
windows. The 1.0- and 2.0-s measurements formed a stable long-window plateau,
whereas 0.25--0.5-s windows gave median peak-to-peak pressures approximately
16\% higher and the 0.125-s window partly followed the impulse. The 1-s moving
median was therefore retained as a reproducible, conservative correction.
Although a centered median is nonlinear and not formally a zero-phase filter,
it does not impose the systematic delay of a one-sided causal smoother.

\subsection{Reduced-time event association}

Before event association, each accepted pick was corrected for its expected differential travel time relative to DD2:
\begin{equation}
  t_{i,\mathrm{red}} = t_i - \frac{r_i-r_{\mathrm{DD2}}}{c_{\mathrm{eff}}},
  \label{eq:reduced_time}
\end{equation}
where $r_i$ is the SLC-40--sensor distance and $c_{\mathrm{eff}}=350.98~\mathrm{m\,s^{-1}}$. The resulting corrections were approximately 0.090~s for DD1, 0~s for DD2, and 0.013~s for DD3. These corrections compensated for the known source--receiver geometry only during preliminary association; final propagation direction and apparent speed were estimated independently from the waveforms.

Reduced picks were sorted chronologically. Beginning with the earliest unassigned pick, all picks falling within the following 0.040~s were treated as a candidate cluster. At most one pick per channel was retained, with the pick closest to the cluster seed selected when multiple picks occurred on the same channel. A cluster was accepted as a candidate event if it contained picks from at least two distinct infrasound channels. Picks that failed this criterion remained unassociated rather than being forced into an event.

Association counts were evaluated for coincidence windows from 0.008 to
0.060~s. The number of multichannel candidates stabilized at 154 for windows
of 0.040--0.060~s, so 0.040~s, the smallest value on that plateau, was adopted.
One candidate was subsequently rejected after waveform review because it had
DD1 and DD2 picks but no coherent DD3 counterpart. The authoritative catalogue
therefore contains 153 events.

For each retained event, a six-channel waveform segment was extracted from
0.50~s before to 0.60~s after the reduced event time. The travel-time
corrections were not written into these waveform streams; observed sample
times were retained. Every segment was required to contain DD1--DD3 and
DHE/DHN/DHZ and to cover the complete requested interval. Pressure stacks used
only the channels assigned to an event, preventing an unassociated DD1 trace
from influencing events supported only by DD2--DD3.

\subsection{Acoustic pulse measurements and event-quality tiers}

For initial pulse characterization, the three infrasound traces were shifted to the DD2 reference time using the meteorologically predicted differential delays and combined using a sample-by-sample median. Within a search interval from $-0.04$ to $+0.25$~s relative to the reduced event time, the algorithm sought a positive compression followed by a negative pressure excursion. The negative peak was required to occur 0.010--0.180~s after the positive peak. When several peak pairs satisfied these criteria, the pair with the largest peak-to-peak pressure was retained.

For positive and negative peaks at times $t_+$ and $t_-$, provisional one-cycle bounds were defined as
\begin{equation}
  t_{\mathrm{start}} = t_+ - 0.5(t_- - t_+)
  \label{eq:pulse_start}
\end{equation}
and
\begin{equation}
  t_{\mathrm{end}} = t_- + 0.5(t_- - t_+).
  \label{eq:pulse_end}
\end{equation}
These model bounds were refined by searching outward for at least three
consecutive samples whose absolute amplitudes were less than twice the robust
background scale, with the search limited to an additional 0.12~s on either
side. Both the model and noise-refined bounds were retained as morphology
diagnostics; the dedicated amplitude analysis remained authoritative for
physical pressure and ground-velocity measurements.

Noise was measured from $-0.40$ to $-0.10$~s relative to the event using the robust scale
\begin{equation}
  \sigma_{\mathrm{robust}} = 1.4826\,\mathrm{median}\left(\left|p-\mathrm{median}(p)\right|\right).
  \label{eq:robust_noise}
\end{equation}
The provisional stack signal-to-noise ratio was the maximum absolute pressure within the event bounds divided by this noise scale.

The positive--negative pulse morphology was used as a measurement-quality criterion rather than as the defining requirement for a physical event. In particular, the principal explosion was retained despite failing the simple one-cycle morphology test because it was an exceptionally strong, coherent, multichannel signal with a more complex pressure history.

Candidate labels based on channel support, reduced-pick span, pressure amplitude, and signal-to-noise ratio were retained as descriptive quality indicators. They were not interpreted as distinct physical source classes.

\subsection{Waveform and array-solution quality}

Catalogue membership was kept separate from later quality classifications.
For waveform analysis, a usable subset required robust peak SNR $\geq3$ and
DD2--DD3 waveform correlation $\geq0.5$; a core waveform subset required SNR
$\geq5$ and correlation $\geq0.7$. DD1 was not used as a mandatory veto
because it contained fewer picks and showed poorer consistency for weak
events.

All 153 final events were submitted to the time-domain planar inversion,
irrespective of amplitude, SNR, or compression--rarefaction morphology. Array
quality was assigned only after inversion from the maximum coherence and its
sensitivity to analysis-window placement. This avoids turning a waveform
morphology or amplitude threshold into a retrospective event-existence
criterion.

\subsection{Continuous Bartlett beamforming}

The catalogue was independently validated using continuous Bartlett
beamforming of DD1--DD3. The primary configuration used 0.5-s windows and a
2--20~Hz band. For each catalogue time, the beamformer estimated back azimuth,
apparent trace velocity, and multichannel coherence $C$. Coherence was
expressed as an array Fisher statistic,
\begin{equation}
  F=\frac{(N-1)C}{1-C},
  \label{eq:f_statistic}
\end{equation}
where $N=3$. The 99.5th percentile of the background $F$ distribution,
$F=3.02$, was used as an empirical reference rather than a catalogue-membership
threshold. Solutions at back-azimuth or velocity-grid limits were flagged.

Sensitivity was tested with 0.25-s windows over 4--30~Hz and 1.0-s windows
over 1--10~Hz. Because the distribution of $F$ depends on window duration and
frequency band, each configuration was interpreted relative to its own
background. The preferred configuration minimized median back-azimuth error,
maximized median $F$, and produced substantially fewer velocity-grid-edge
solutions. Local maxima in the continuous beamformer were matched one-to-one
with catalogue times as a completeness test, but were not counted as physical
explosions because overlapping windows can generate multiple maxima from one
oscillatory event.

\subsection{Planar-wave array analysis}

Back azimuth and apparent horizontal propagation speed were estimated jointly by shifting all three infrasound waveforms according to trial plane-wave models. Back azimuth $\theta$ was defined as the direction from the array toward the source, measured clockwise from north. For a trial back azimuth and apparent speed $v$, the horizontal propagation vector toward the array was
\begin{equation}
  \mathbf{k} = [-\sin\theta,-\cos\theta],
  \label{eq:propagation_vector}
\end{equation}
and the predicted relative delay at sensor $i$, relative to DD2, was
\begin{equation}
  \Delta t_i = \frac{(\mathbf{x}_i-\mathbf{x}_{\mathrm{DD2}})\cdot\mathbf{k}}{v}.
  \label{eq:plane_delay}
\end{equation}

For each trial model, the three waveforms were shifted by interpolation and trimmed to their common temporal overlap. Every point on the search surface was therefore evaluated using the same waveform duration, avoiding the unequal-overlap bias that can favor zero lag in conventional cross-correlation calculations.

Each aligned trace was median-centered and normalized by its root-mean-square amplitude. Coherence was evaluated using both three-channel semblance,
\begin{equation}
  S = \frac{\sum_t\left[\sum_i x_i(t)\right]^2}{N\sum_t\sum_i x_i^2(t)},
  \label{eq:semblance}
\end{equation}
and the mean of the three pairwise correlation coefficients. The mean correlation $\overline{\rho}$ was mapped from $[-1,1]$ to $[0,1]$, and the combined score was
\begin{equation}
  C = 0.60S + 0.40\left(\frac{\overline{\rho}+1}{2}\right).
  \label{eq:coherence_score}
\end{equation}

The coarse search covered back azimuths from 180$^\circ$ to 220$^\circ$ in 0.5$^\circ$ increments and apparent speeds from 280 to 440~m~s$^{-1}$ in 2~m~s$^{-1}$ increments. Each coarse maximum was refined over $\pm2^\circ$ in 0.1$^\circ$ increments and $\pm20$~m~s$^{-1}$ in 0.5~m~s$^{-1}$ increments. Most events were scored over the interval $-0.04$ to $+0.18$~s relative to the reduced event time; the principal explosion used an extended interval ending at $+0.35$~s.

Pairwise correlations and residual time lags were calculated after application of the best-fitting model. The lag-closure quantity
\begin{equation}
  \tau_{12}+\tau_{23}+\tau_{31}
  \label{eq:lag_closure}
\end{equation}
was retained as a consistency diagnostic. The relative performance of DD1 was evaluated by comparing the DD2--DD3 correlation with the mean correlation of the two pairs containing DD1.

Solution stability was tested by shifting the scoring window by $-0.008$, 0, and $+0.008$~s and scaling its duration by factors of 0.9, 1.0, and 1.1. Standard deviations of the resulting back azimuths and speeds were retained as empirical stability measures. Widths of the coherence surface at score reductions of 0.01 and 0.02 were also calculated, but these widths were treated as descriptive measures rather than formal confidence intervals.

A high-quality solution required $C\geq0.90$, window-sensitivity back-azimuth
standard deviation $\leq0.5^\circ$, and velocity standard deviation
$\leq3$~m~s$^{-1}$. An intermediate solution required $C\geq0.75$,
back-azimuth variability $\leq1^\circ$, and velocity variability
$\leq8$~m~s$^{-1}$. All remaining solutions were retained as weak or broad.
Continuous coherence and stability values were preserved. Catalogue-level
means used maximum coherence as an event weight; the associated weighted
standard deviations describe event-to-event dispersion, not uncertainty on the
weighted mean.

\subsection{Finite-range circular-wavefront analysis}

The adequacy of the plane-wave approximation was assessed using exact horizontal source--receiver distances. For a trial source position $\mathbf{x}_s$ and propagation speed $v$, the predicted delay at sensor $i$, relative to DD2, was
\begin{equation}
  \Delta t_i = \frac{|\mathbf{x}_i-\mathbf{x}_s|-|\mathbf{x}_{\mathrm{DD2}}-\mathbf{x}_s|}{v}.
  \label{eq:circular_delay}
\end{equation}
Because receiver elevations were not included in the adopted geometry, this was a two-dimensional circular-wavefront model rather than a full three-dimensional spherical-wave model.

The primary finite-range calculation fixed the source at SLC-40 and searched propagation speeds from 280 to 440~m~s$^{-1}$ in 0.5~m~s$^{-1}$ increments. Coherence was evaluated using the same waveform windows, normalization, and scoring metric as the planar-wave analysis.

Sensitivity to source position was examined in coordinates oriented parallel and perpendicular to the SLC-40--BCHH path. Along-range offsets from $-500$ to $+500$~m and cross-range offsets from $-150$ to $+150$~m were evaluated while profiling over the complementary coordinate. These calculations were interpreted as resolution tests rather than independent source locations. With three sensors, only two independent relative delays are available, so source easting, source northing, and propagation speed cannot all be determined independently. Open along-range coherence ridges were therefore interpreted as a fundamental range-resolution limitation of the array geometry.

\subsection{Acoustic and seismic amplitude measurements}

Acoustic and seismic amplitudes were measured for each catalogue event using
the saved baseline-corrected waveform segments. Event intervals were taken from
the positive--negative pulse measurements, except for the principal explosion,
for which a fixed interval from $-0.04$ to $+0.35$~s was used.

For each seismic component, the median ground velocity from $-0.18$ to $-0.05$~s was removed. Component peak ground velocity was the largest absolute velocity within the event interval. Three-component vector peak ground velocity was calculated as
\begin{equation}
  \mathrm{PGV}_{\mathrm{vec}} = \max_t\sqrt{v_E^2(t)+v_N^2(t)+v_Z^2(t)}.
  \label{eq:vector_pgv}
\end{equation}
The maximum vector amplitude over the complete saved waveform segment was also retained to determine whether the acoustically defined event interval captured the largest seismic motion.

For catalogue measurements, the three infrasound traces were aligned to DD2
using the empirical array-scale apparent velocity of approximately
352~m~s$^{-1}$ and combined by a sample-by-sample median. Positive pressure,
negative pressure, and peak-to-peak pressure were measured within the same
event interval. Acoustic signal-to-noise ratio was defined as median-stack
peak-to-peak pressure divided by the robust pre-event noise scale.

The median-stack measurements provide a consistent catalogue metric but differ
from measurements made separately on the three physical sensors. For the
principal explosion and the other named events, absolute pressures were
therefore characterized from the individual calibrated sensors and summarized
by their median and range. The final principal-explosion medians were
+1.392~kPa, $-157$~Pa, and 1.549~kPa peak-to-peak; positive and peak-to-peak
ranges were 1.31--1.51 and 1.44--1.67~kPa, respectively. Sensor-specific and
median-stack quantities represent different measurement procedures and were
not used interchangeably.

Where pressure was reduced to a reference distance of 1~km, spherical geometric spreading was assumed:
\begin{equation}
  p_{\mathrm{red},1\,\mathrm{km}} = p_{\mathrm{obs}}\frac{r}{1000~\mathrm{m}}.
  \label{eq:reduced_pressure}
\end{equation}
For the median principal-explosion positive peak, this gives a reduced positive
pressure of approximately 1.967~kPa at 1~km.

Peak pressure level was calculated as
\begin{equation}
  L_{\mathrm{peak}} = 20\log_{10}\left(\frac{p_+}{20~\mu\mathrm{Pa}}\right).
  \label{eq:peak_pressure_level}
\end{equation}
This is a peak pressure level, not an RMS or time-averaged sound-pressure level. The positive peak pressure of approximately 1435~Pa corresponds to a peak level of approximately 157~dB re 20~$\mu$Pa.

A dimensional acoustic-to-seismic amplitude ratio was calculated as peak-to-peak pressure divided by vector peak ground velocity, with units of Pa per m~s$^{-1}$. This quantity was treated as an empirical measure of pressure--ground-motion coupling and not as an acoustic-to-seismic energy ratio.

\subsection{Acoustic--seismic amplitude regression}

Regression analysis was restricted to events with vector-peak-ground-velocity signal-to-noise ratio of at least 5, acoustic peak-to-peak signal-to-noise ratio of at least 5, and an event-window peak ground velocity equal to at least 70\% of the maximum peak ground velocity measured in the complete saved segment.

A free power-law relationship,
\begin{equation}
  p_{\mathrm{p2p}} = A\,\mathrm{PGV}^{b},
  \label{eq:power_law}
\end{equation}
was fitted by ordinary least squares after transforming both variables to base-10 logarithms. A 95\% descriptive interval for the exponent $b$ was estimated using 5000 nonparametric bootstrap resamples of the event population.

A proportional model,
\begin{equation}
  p_{\mathrm{p2p}} = K\,\mathrm{PGV},
  \label{eq:proportional_model}
\end{equation}
was also evaluated as a constant-ratio alternative. The two models were compared using root-mean-square residuals in base-10 logarithmic space. These regressions describe the observed relationship only; they do not account for measurement uncertainty in both variables, covariance among sensor-calibration errors, or censoring near the acoustic or seismic noise floors.

\subsection{Spectra and empirical pressure-to-ground transfer estimates}

Pressure and vertical ground-velocity amplitude spectra were calculated for the
51 high-quality array events. Records were linearly detrended, tapered with a
Tukey window ($\alpha=0.20$), and transformed using periodograms. Spectral
amplitudes from 0.9 to 90~Hz were summarized in logarithmic bins using the
geometric median and 16th--84th percentiles. The approximately 1.10-s event
records gave an independent frequency spacing of approximately 0.91~Hz.

Pressure-to-vertical-velocity coupling was estimated using a linearly averaged,
delay-aligned pressure beam as input. The H1 estimator and magnitude-squared
coherence were
\begin{equation}
  H_1(f)=\frac{S_{pv}(f)}{S_{pp}(f)}, \qquad
  \gamma^2(f)=\frac{|S_{pv}(f)|^2}{S_{pp}(f)S_{vv}(f)}.
  \label{eq:h1_transfer}
\end{equation}
Auto- and cross-spectra were calculated for each event and averaged before the
transfer estimate was formed. Event bootstrap resampling supplied 95\%
intervals. Frequency-specific coherence thresholds were obtained from 1000
permutations that paired pressure and velocity from different events.

A complementary continuous estimate used the 70-s Phase~I record divided into
seven nonoverlapping 10-s segments. These records were linearly detrended
without reapplying the infrasound baseline correction. Segment bootstrap and
permutation tests paralleled the event analysis. The event and continuous
estimates overlap in time and were treated as complementary analyses, not
independent replication. Sensitivity tests excluded the strongest one, three,
and five events, equalized pressure RMS, and substituted a median pressure
stack for the linear beam.

\subsection{Sound exposure, acoustic energy, and yield-equivalent scenarios}

Sound exposure and radiated acoustic energy were estimated for the union of
high-quality array events and the five named events. Because adjacent catalogue
windows can sample the same physical blast complex, selected DD2-time windows
were first merged into nonoverlapping intervals. The pressure channels were
differentially aligned and $p^2(t)$ was integrated once per event complex on
each channel. The hemispherical estimate was
\begin{equation}
  E_{2\pi}=2\pi r^2\int\frac{p^2(t)}{\rho c}\,dt,
  \label{eq:acoustic_energy}
\end{equation}
where $\rho$ is moist-air density and $c=351$~m~s$^{-1}$. Medians and ranges
across DD1--DD3 were retained. Cumulative energy was summed channel by channel
over disjoint complexes, preventing the principal explosion from being counted
more than once.

Radiated acoustic-energy TNT equivalent was calculated using
$4.184\times10^6$~J~kg$^{-1}$ TNT. For sensitivity only, total explosive yield
was estimated by dividing this acoustic-energy equivalent by efficiencies of
0.1, 1, and 10\%. These are transparent spreading and efficiency scenarios,
not calibrated blast yields. The earlier MATLAB estimates based on
\citet{johnson_aster_2005} and \citet{kanamori_1977} are retained only for
historical comparison and are not combined with the Python event-complex
energetics.

\subsection{Video and audio comparison}

Publicly available video recordings were used to identify major visible stages
of the accident. The USLaunchReport video lacks an absolute UTC time base, so
its first upper-stage flash was mapped to the geophysical source chronology
rather than treated as an independent clock. Identifiable changes included the
first flash, expansion of the principal fireball, structural collapse, descent
of vehicle components, apparent ground contact, and subsequent explosions.

The quantitative event catalogue, pressure measurements, and array results were
derived independently from the BCHH geophysical data. Video observations were
used principally to interpret the physical processes associated with key
signals and to organize the accident chronology. Camera audio was advanced by
its own path travel time, calculated from a source-to-camera distance of
approximately 4.255~km and a path-specific meteorological effective speed of
349.8~m~s$^{-1}$, corresponding to approximately 12.16~s. This differs from
the BCHH correction because the two paths have different wind projections.
Once the synchronized source-aligned audio product had been generated, it was
not shifted again in publication figures.

Source-aligned comparisons shifted each BCHH trace by its channel-specific
distance divided by 351~m~s$^{-1}$. Camera audio was high-pass filtered above
300~Hz for display, pressure was shown without additional filtering beyond
baseline correction, and seismic traces were filtered from 0.5 to 20~Hz.
Compact event panels show the video frame immediately before and after each
mapped event together with source-reduced audio, pressure, and three-component
ground velocity (Figure~\ref{fig:multimodal}). Traces were independently
normalized, so these panels demonstrate timing rather than relative amplitude.

\subsection{Polarization analysis and synchronized forensic visualization}
\label{sec:methods_polarization}

Response-corrected ground velocity was rotated into vertical, radial, and
transverse coordinates using the 199.438$^\circ$ back azimuth from BCHH toward
SLC-40. Sequence-wide polarization was evaluated after 0.5--20~Hz filtering
using overlapping real covariance and complex analytic-signal windows.
Frequency dependence was examined in 0.5--3, 3--10, 10--25, 25--45, and
0.5--20~Hz bands. Eigenvalues and eigenvectors were used to calculate
linearity, rectilinearity, planarity, principal azimuth, and incidence angle;
the analytic-signal calculation provided degree of polarization and
ellipticity. Low-amplitude windows were retained for overview displays but were
not interpreted as event polarization.

A focused search examined the interval between the upper-stage source time and
its predicted ground-coupled-airwave arrival 4.046~s later. Because the much
larger principal explosion occurred at +3.601~s, the search ended at +3.45~s.
Causal two-pole Butterworth filters were applied in 0.8--3, 3--8, 8--20, and
20--45~Hz bands; zero-phase results were retained as timing-sensitivity tests.
Persistent vector-amplitude crossings at four, five, and six robust background
standard deviations were compared with STA/LTA and AIC estimates. Trailing
polarization windows prevented an estimate from using samples later than its
reported time. Candidate motion required amplitude above four background
standard deviations, rectilinearity $\geq0.70$, degree of polarization
$\geq0.55$, and radiality $\geq0.70$. Apparent velocities were calculated from
the 1419.9-m path length. Phase names were not assigned from the single-station
observations.

\subsection{Interpretation and provenance}

Throughout the analysis, direct observations were distinguished from later physical interpretations. Direct observations included arrival times, waveform amplitudes, pressure polarity and duration, apparent propagation speeds, back azimuths, coherence measures, ground velocities, and derived energy estimates. Interpretations of individual signals were based on their temporal correspondence with publicly available video and, where relevant, information released during the subsequent SpaceX failure investigation.

The modern Python workflow was treated as authoritative for response
correction, baseline removal, event association, waveform review, Bartlett and
time-domain array analyses, finite-distance sensitivity, acoustic and seismic
amplitudes, amplitude regression, spectra, empirical transfer estimates,
acoustic energetics, polarization, and multimedia alignment. Legacy MATLAB
values were retained only for documented historical comparison. Where the two
workflows produced related quantities, their definitions and provenance were
kept explicit rather than silently combined.

\subsection{Reproducibility}

Each Python processing stage writes its numerical products, configuration parameters, intermediate event streams, event-selection decisions, and diagnostic figures to a dedicated output directory. The canonical workflow comprises notebooks for response correction, input preparation, meteorological processing, catalogue reconciliation and validation, waveform review, baseline validation, fixed-event products, planar- and finite-distance array analysis, acoustic--seismic amplitude and spectral analysis, key-event measurements, polarization, synchronized video, and final paper figures. Quality classifications and diagnostic metrics are retained for every catalogue event.

The archive will also include the Python implementation of the
three-component polarization analysis, the resulting polarization attributes
with complete processing metadata, selected particle-motion products, and the
script used to generate the synchronized supplementary video. The source video identifier, synchronization
offset, frame-rate handling, waveform time base, filtering, moving-window
parameters, coordinate conventions, and definitions of all polarization
attributes will be documented explicitly. Where redistribution of the original
public video is not permitted, the archive will provide the synchronization
metadata and rendering scripts together with an exact citation and identifier
for the source video rather than redistributing the source footage.

The continuous waveforms, station metadata, manual arrival picks, event catalogue, Python notebooks, legacy MATLAB routines, derived tables, and figure-generation scripts will be archived in an open repository with a persistent identifier. The archive will identify the specific workflow and software implementation used to generate each principal result. The Python processing uses ObsPy \citep{beyreuther_obspy_2010,krischer_obspy_2015}
and the scientific Python ecosystem, whereas the original analysis used
Antelope \citep{brtt_antelope}, GISMO \citep{thompson_gismo_2017}, and custom
MATLAB routines. Final software and data citations will also document NumPy,
SciPy, pandas, Matplotlib, pyproj, the KSC meteorological observations, the
instrument responses and infraBSU calibration information, and the public
videos used in the chronology.

\section{Results}
\label{sec:results}

\subsection{Atmospheric conditions}
\label{sec:results_atmosphere}

Atmospheric conditions remained nearly constant throughout the 30-min
sequence. Mean temperature was 26.36$\pm$0.27~$^\circ$C and mean relative
humidity was 83.9$\pm$3.7\%. The lower-atmosphere profile gave a vector-mean
wind of 4.57~m~s$^{-1}$ from 144.9$^\circ$, whose projection onto the
SLC-40--BCHH path was a +2.65~m~s$^{-1}$ tailwind. These conditions correspond
to a still-air sound speed of 348.33~m~s$^{-1}$ and an effective speed of
350.98~m~s$^{-1}$. Over 1419.9~m, the wind reduced travel time from 4.076 to
4.046~s. Little systematic temporal variation was observed, supporting use of
one atmospheric model. The prediction is effectively identical to the rounded
351~m~s$^{-1}$ reduction speed and agrees closely with the independent
array-scale apparent velocity of approximately 352~m~s$^{-1}$.


\subsection{Final event catalogue and temporal evolution}
\label{sec:results_catalogue}

Manual review yielded 384 accepted infrasound picks. After reduced-time
association and waveform review, the final catalogue contained 153 events
supported by 369 picks: 90 events were represented on DD2 and DD3, and 63 on
all three pressure channels. No retained event was supported only by DD1 and
DD2. The rejected candidate lacked a coherent DD3 counterpart and had a
DD2--DD3 waveform correlation of only 0.092.

The activity was strongly episodic (Figure~\ref{fig:overview}). Phase~I
contained 48 events, Phase~II 4, Phase~III 96, and Phase~IV 5. The opening
sequence contained the largest signals, whereas Phase~III contained the
highest event rate. Sustained seismic tremor occurred between approximately
13:08:48 and 13:12:00~UTC without a corresponding sequence of impulsive
infrasound arrivals.

Of the 153 events, 138 met the usable-waveform criteria and 107 met the core
waveform criteria. Median robust peak SNR was 7.86, median DD2--DD3 correlation
was 0.876, and median residual DD2--DD3 lag after geometric correction was
zero. The median corrected span among associated picks was 4.87~ms
(interquartile range 2.13--8.87~ms; maximum 36.24~ms), providing an internal
check on both event association and differential travel-time reduction. A
simple compression--rarefaction morphology was measurable for 152 events; the
principal explosion was the sole complex or noncanonical waveform.



\begin{figure*}[t]
  \centering
  \includegraphics[width=0.98\textwidth]{fig03_bchh_30min_overview_zp_true.pdf}
  \caption{Overview of the Falcon~9 accident sequence recorded at BCHH. The
  upper three traces show acoustic pressure on DD1--DD3 and the lower three
  traces show three-component ground velocity. The four displayed phase
  intervals and the prolonged seismic-tremor interval are indicated. Phases
  I--IV contain 48, 4, 96, and 5 catalogue events, respectively. Time is UTC
  on 1 September 2016.}
  \label{fig:overview}
\end{figure*}

\subsection{Absence of detectable precursors and later activity}
\label{sec:results_onset}

To test whether a detectable energetic precursor preceded the visible failure,
the complete continuous dataset spanning approximately 37~h before the first
accident-related signal was reviewed for anomalous impulsive seismic or
infrasonic activity. No such precursor attributable to the accident was
identified. The pre-event record contained normal background seismicity,
cultural noise, and environmental signals, but no coherent impulsive arrival
from the direction of SLC-40 that could be associated with an external
energetic event.

This result places an observational constraint on detectable seismo-acoustic
precursors at BCHH. It does not identify the initiating engineering mechanism
and cannot exclude aseismic, non-impulsive, weak, or otherwise undetectable
processes within the launch vehicle.

No additional explosion-related impulsive arrivals were identified during the
10.5~h following the final event in Phase~4.


\subsection{Opening sequence and key accident stages}
\label{sec:results_principal}

The opening sequence contained five named signals: the initial upper-stage
event, the principal lower-stage explosion, capsule impact, capsule explosion,
and event~024 following the capsule explosion. Their source times were
13:07:11.913, 13:07:15.514, 13:07:24.420, 13:07:24.988, and 13:07:26.458~UTC,
respectively (Table~\ref{tab:principal_events}; Figures~\ref{fig:opening}
and~\ref{fig:multimodal}). These times were inferred from precisely reviewed
BCHH arrivals using the 351~m~s$^{-1}$ source-to-DD2 reduction; the video was
subsequently mapped to this chronology.

The upper-stage event produced median channel-specific positive, negative, and
peak-to-peak pressures of +38.8, $-25.4$, and 77.1~Pa, respectively. Its
three-component vector PGV was
$7.85\times10^{-5}\ \si{\metre\per\second}$. Geometric reduction gave a
median peak-to-peak pressure of 108.6~Pa at 1~km.

Approximately \SI{3.60}{\second} later, the principal explosion produced the
largest acoustic and seismic amplitudes in the catalogue
(Figure~\ref{fig:principal}). Median channel-specific positive, negative, and
peak-to-peak pressures were \SI{1.392}{\kilo\pascal},
\SI{-157}{\pascal}, and \SI{1.549}{\kilo\pascal}, respectively. Across
DD1--DD3, positive pressure ranged from 1.31 to 1.51~kPa and peak-to-peak
pressure from 1.44 to 1.67~kPa. Assuming $1/r$ geometric spreading, the median
positive and peak-to-peak pressures were 1.967 and 2.189~kPa at 1~km. The
positive peak corresponds to approximately 157~dB re 20~$\mu$Pa.

The principal pressure waveform comprised an abrupt positive compression
followed by a longer, irregular decay and substantially smaller rarefaction. It
therefore did not display the compact, approximately symmetric positive--negative
morphology of many weaker catalogue events. Its shape qualitatively resembles a
blast pulse with an extended decay, although no parametric Friedlander-waveform
fit was performed. The seismic response persisted considerably longer than the
initial acoustic compression and contained large-amplitude motion on all three
components. The three-component vector PGV reached
$3.57\times10^{-3}\ \si{\metre\per\second}$, the largest value measured in the
catalogue.

The two capsule-related signals were separated by approximately
\SI{0.568}{\second}. Capsule impact produced pressure amplitudes of order
$+276$, $-102$, and 378~Pa peak-to-peak and vector PGV of
$6.98\times10^{-4}\ \si{\metre\per\second}$. Capsule explosion produced
median positive, negative, and peak-to-peak pressures of +297.0, $-113.4$,
and 410.6~Pa and vector PGV of
$9.02\times10^{-4}\ \si{\metre\per\second}$. Both events had coherent
three-sensor waveforms and apparent velocities close to the 351--352~m~s$^{-1}$
meteorological and empirical values, consistent with propagation from SLC-40
(Figure~\ref{fig:capsule}). Event~024 formed a third distinct
post-capsule pressure pulse but was omitted from the uncluttered 30-s overview
annotation.



\begin{figure*}[t]
  \centering
  \includegraphics[width=0.98\textwidth]{fig03_bchh_30s_zoom_with_audio_reduced_zp_true.pdf}
  \caption{Source-aligned 30-s view of the opening accident sequence. The
  synchronized USLaunchReport audio is shown above DD1--DD3 pressure and
  three-component ground velocity. BCHH traces were reduced using
  receiver-specific distances and 351~m~s$^{-1}$; camera audio had previously
  been advanced by its approximately 12.16-s travel time using the
  camera-path effective speed of 349.8~m~s$^{-1}$. Markers identify only the
  upper-stage event, principal explosion, capsule impact, and capsule
  explosion. Because the video lacks an independent absolute timestamp, its
  mapped UTC scale is conditional on the upper-stage visual anchor.}
  \label{fig:opening}
\end{figure*}

\begin{figure*}[t]
  \centering
  \includegraphics[width=0.98\textwidth]{fig04_multimodal_alignment_upper_stage.pdf}
  \caption{Video, camera audio, infrasound, and seismic alignment for the
  upper-stage event. Two native USLaunchReport frames immediately before and
  after the mapped event show the first visible flash. Camera audio was
  advanced by its approximately 12.16-s path travel time using the 349.8
  m~s$^{-1}$ camera-path effective speed; BCHH records were reduced using
  channel-specific distances and 351~m~s$^{-1}$. Camera audio is high-pass
  filtered above 300~Hz, pressure is unfiltered apart from baseline correction,
  and seismic traces are filtered from 0.5 to 20~Hz. Traces are independently
  normalized, so the panel demonstrates timing rather than relative amplitude.}
  \label{fig:multimodal}
\end{figure*}


\begin{figure*}[t]
  \centering
  \includegraphics[width=0.92\textwidth]{principal_explosion_waveforms.png}
  \caption{Calibrated pressure and three-component ground-velocity waveforms
  for the principal explosion. Median channel-specific positive, negative, and
  peak-to-peak pressures are +1.392~kPa, $-157$~Pa, and 1.549~kPa,
  respectively; DD1--DD3 ranges show the array-scale measurement variability.}
  \label{fig:principal}
\end{figure*}


\begin{figure*}[t]
  \centering
  \includegraphics[width=0.92\textwidth]{capsule_explosion_waveforms.png}
  \caption{Calibrated waveforms for capsule impact and capsule explosion,
  separated by 0.568~s. Their pressure amplitudes are summarized from the
  three physical sensors, and their vector PGVs are
  $6.98\times10^{-4}$ and $9.02\times10^{-4}$~m~s$^{-1}$, respectively. The
  physical labels are based on the mapped video chronology.}
  \label{fig:capsule}
\end{figure*}


\begin{table*}[t]
\caption{Named signals in the opening sequence. Pressure values are medians of
the three channel-specific scalar measurements; PGV is the three-component
vector peak ground velocity. Source times were obtained from reviewed BCHH
arrivals using the 351~m~s$^{-1}$ source-to-DD2 reduction. Capsule-impact
pressure remains sensitive to the leading edge of its manual signal window.
Event~024 was included in catalogue and energetics analyses but was not given a
separate manual pressure window in the key-event measurement notebook.}
\label{tab:principal_events}
\centering
\small
\begin{tabular}{lccccc}
\hline
Signal & Source time (UTC) & $p_+$ (Pa) & $p_-$ (Pa) & $p_{\mathrm{p2p}}$ (Pa) & Vector PGV (m s$^{-1}$) \\
\hline
Upper-stage event & 13:07:11.913 & 38.8 & $-25.4$ & 77.1 & $7.85\times10^{-5}$ \\
Principal explosion & 13:07:15.514 & 1392 & $-157$ & 1549 & $3.57\times10^{-3}$ \\
Capsule impact & 13:07:24.420 & $\sim276$ & $\sim-102$ & $\sim378$ & $6.98\times10^{-4}$ \\
Capsule explosion & 13:07:24.988 & 297.0 & $-113.4$ & 410.6 & $9.02\times10^{-4}$ \\
Event 024 & 13:07:26.458 & --- & --- & --- & --- \\
\hline
\end{tabular}
\end{table*}

\subsection{Seismo-acoustic waveform characteristics}
\label{sec:results_waveforms}

Many stronger impulsive events were recorded simultaneously by the pressure
array and broadband seismometer, illustrating efficient air-to-ground coupling.
The median of 51 independently normalized high-quality waveforms contained a
compact positive pressure pulse beginning near reduced time zero and peaking
approximately 0.02~s later, followed by a smaller negative excursion. Median
vertical ground velocity began near the pressure onset but remained oscillatory
for at least 0.3~s, with characteristic periods of approximately 0.04--0.08~s.
Persistence of this coda in the normalized ensemble demonstrates repeatable
waveform structure rather than domination by the principal explosion.

\begin{figure*}[t]
  \centering
  \includegraphics[width=0.88\textwidth]{high_quality_normalized_pressure_and_dhz_waveforms.pdf}
  \caption{Typical waveforms for the 51 high-quality planar-array events.
  Pressure and vertical ground velocity were normalized independently within
  each event before stacking, so the panels emphasize waveform shape rather
  than absolute amplitude. Pale curves show individual events, the dark curve
  is the pointwise median, and shading spans the 16th--84th percentiles. The
  compact pressure pulse is followed by a longer, oscillatory vertical ground
  response. Time is relative to the DD2-referenced reduced event time.}
  \label{fig:waveform_ensemble}
\end{figure*}

The principal explosion produced the strongest seismic response, including
high-amplitude surface-wave-like motion that dominated ground velocity for
approximately 20~s. In contrast, many smaller impulsive events produced clear
infrasound arrivals but only weak or undetectable seismic signals
(Figure~\ref{fig:waveform_ensemble}).

Public video shows the payload or capsule descending and apparently reaching
the ground during the two capsule-associated pulses. The raw three-component
waveforms remain dominated by ground-coupled airwaves, so the physical label
``capsule impact'' follows the mapped visual chronology rather than a uniquely
isolated seismic impact phase.

Following the initial phase of impulsive activity, a prolonged continuous
signal persisted for approximately 3.5~min. It is strongest on the seismic
channels but may also be present on DD2 and therefore should not presently be
classified as seismic-only. Preliminary comparison with video suggests a
possible association with sustained launch-pad fire was considered, but its
physical cause is not uniquely resolved by the present data.

The infrasound waveforms were predominantly impulsive, with typical rise times of approximately 0.01--0.02~s. Even for the largest explosion, however, the waveforms did not exhibit the classical symmetric N-wave morphology commonly associated with idealized blast waves.


\subsection{Polarization and particle-motion results}
\label{sec:results_polarization}

The source-oriented seismic records contained organized motion before the
upper-stage ground-coupled airwave was expected at +4.046~s
(Figure~\ref{fig:preairwave}). Persistent causal amplitude increases occurred
at 1.843--1.947~s in the 0.8--3~Hz band and 1.855--1.999~s in the 3--8~Hz
band, with medians of 1.863 and 1.895~s. Zero-phase estimates were 1.739 and
1.703~s. Across filter and detector choices, a conservative onset interval is
therefore 1.7--2.0~s after the upper-stage source time.

At 1419.9~m range, this interval corresponds to an apparent velocity of
approximately 0.71--0.84~km~s$^{-1}$, faster than the atmospheric airwave but
slower than typical competent-rock body waves. Amplitude-gated motion was
predominantly rectilinear and radial-horizontal, followed by approximately
quadrature Z--R motion. These observations favor a shallow S/SV or
Rayleigh-type packet, or mixed near-surface motion, rather than an unambiguous
direct P arrival. A later 8--20~Hz onset at approximately 2.53--2.79~s
corresponded to 0.51--0.56~km~s$^{-1}$. The 20--45~Hz result was not
repeatable and was not interpreted.

\begin{figure*}[t]
  \centering
  \includegraphics[width=0.96\textwidth]{12_multiband_polarization_comparison.png}
  \caption{Source-oriented search for seismic motion preceding the predicted
  upper-stage ground-coupled airwave. Noise-normalized vertical, radial, and
  transverse measures are shown in multiple frequency bands relative to the
  upper-stage source time. The conservative 1.7--2.0~s onset interval precedes
  the predicted 4.046-s airwave and corresponds to an apparent velocity of
  approximately 0.71--0.84~km~s$^{-1}$. Predominantly radial-plane motion
  supports a shallow or surface-wave packet, but no unique phase is assigned.}
  \label{fig:preairwave}
\end{figure*}


\subsection{Planar-wave array results}
\label{sec:results_array}

The independent Bartlett analysis showed that the catalogue was dominated by
signals arriving from SLC-40 (Figure~\ref{fig:array_summary}). Under the
preferred 0.5-s, 2--20~Hz configuration, median back azimuth was
200.67$^\circ$, compared with the geometric 199.44$^\circ$. Of 153 events,
145 (94.8\%) were within 5$^\circ$ of SLC-40 and 151 (98.7\%) within
10$^\circ$. Median apparent velocity was 354.6~m~s$^{-1}$ and median $F$ was
5.04. In total, 113 events exceeded the 99.5th-percentile background threshold
of $F=3.02$, 77 had $F\geq5$, and 42 had $F\geq10$. Apparent-velocity scatter
decreased from 30.7~m~s$^{-1}$ for all events to 12.1~m~s$^{-1}$ for
$F\geq10$ events.

The time-domain inversion returned solutions for all 153 events: 51 high
quality, 45 intermediate, and 57 weak or broad. Catalogue-wide median back
azimuth was 200.5$^\circ$ and the coherence-weighted circular mean was
200.44$^\circ$, with 1.51$^\circ$ event-to-event dispersion. The high-quality
events gave 200.73$^\circ$ with 0.70$^\circ$ dispersion. Their
coherence-weighted apparent velocity was 351.98~m~s$^{-1}$, with
3.85~m~s$^{-1}$ dispersion and a median of 352.5~m~s$^{-1}$. This agrees to
within approximately 1~m~s$^{-1}$ with the 350.98~m~s$^{-1}$ meteorological
prediction. Three weak solutions reached velocity-grid limits and were not
interpreted physically.

The principal explosion produced the most coherent solution: back azimuth
200.3$^\circ$, apparent velocity 352.5~m~s$^{-1}$, and coherence score 0.9984.
Pairwise correlations were 0.9968--0.9978. Predicted DD1 and DD3 delays relative
to DD2 were 0.08845 and 0.01115~s, close to the geometry--weather values of
0.09011 and 0.01287~s. Perturbing the scoring interval produced standard
deviations of 0.082$^\circ$ and 0.41~m~s$^{-1}$; the broad near-maximum surface
shows that these quantify repeatability rather than formal spatial resolution.

The approximately 1$^\circ$ clockwise offset from the nominal pad direction
corresponds to approximately 25~m at 1.4~km range, but only about 0.5~m of
transverse projection across the 30-m aperture. The offset is therefore readily
explained by the 1--2~m uncertainty in reconstructed sensor positions and is not
evidence for a source systematically displaced from SLC-40.

\begin{figure*}[t]
  \centering
  \includegraphics[width=0.98\textwidth]{planar_catalogue_summary.png}
  \caption{Direction, apparent velocity, and coherence of the 153-event
  catalogue. (a) Time-domain planar-wave back azimuth; the dashed line is the
  geometric direction to SLC-40. (b) Apparent velocity; the dashed line is the
  meteorological effective speed of 350.98~m~s$^{-1}$. (c) Bartlett $F$
  statistic; the dashed line is the 99.5th percentile of the background
  distribution ($F=3.02$). Dark symbols denote high-quality solutions,
  intermediate tones denote intermediate solutions, and pale symbols denote
  weak or broad solutions that remain catalogue members. Open squares identify
  velocity-grid-edge solutions. Shading marks Phases I--IV.}
  \label{fig:array_summary}
\end{figure*}


\subsection{Finite-range circular-wavefront results}
\label{sec:results_finite_range}

Replacing the plane-wave approximation with exact horizontal distances to a
source fixed at SLC-40 produced only modest changes. The median preferred
finite-distance speed was 355.5~m~s$^{-1}$, compared with 352.5~m~s$^{-1}$
for the planar model; the median eventwise difference was +4.0~m~s$^{-1}$.
Finite-distance coherence was lower than planar coherence by a median of
0.00512. This is expected because the planar model can vary both direction and
speed, whereas the fixed-source calculation imposes the nominal pad direction.

The modest speed increase is not evidence that the effective atmospheric speed
was approximately 356~m~s$^{-1}$; speed can compensate for the approximately
1$^\circ$ difference between the measured wavefront direction and the direction
imposed by the nominal coordinates. The planar estimate, which independently
matches the meteorological prediction, remains the preferred propagation
result.

The range-profile analysis demonstrated that the array could not determine
source range. At a score reduction of 0.01, every one of the 153 along-range
profiles remained open at the search boundary. Cross-range position was better
constrained, with a median preferred offset of +25~m, but this is equivalent to
the same approximately 1$^\circ$ geometry bias described above. The
finite-distance analysis therefore supports the simpler planar approximation
and demonstrates that BCHH constrains source direction much more effectively
than source range (Figure~\ref{fig:finite_distance}).

\begin{figure*}[t]
  \centering
  \includegraphics[width=0.96\textwidth]{supplement_planar_vs_fixed_source_and_range_sensitivity.png}
  \caption{Sensitivity of the array solutions to finite source distance.
  (a) Preferred planar apparent velocity compared with propagation speed for a
  source fixed at SLC-40; the dashed line denotes equality. (b) Corresponding
  planar and fixed-source coherence scores. (c) Median event-relative coherence
  deficit, $C-C_{\max}$, for trial source positions expressed as cross-range and
  along-range offsets from SLC-40. The star marks the nominal pad coordinate and
  the contour encloses positions within 0.01 of the median maximum. The open
  along-range ridge demonstrates that source range is unresolved. This figure
  is intended for the Supplementary Material.}
  \label{fig:finite_distance}
\end{figure*}


\subsection{Event amplitudes through time}
\label{sec:results_temporal}

Across all 153 catalogue events, median stack peak-to-peak pressure was
\SI{29.5}{\pascal}, with an interquartile range of
\SIrange{16.2}{51.3}{\pascal}. Median vector PGV was
$5.6\times10^{-5}\ \si{\metre\per\second}$. Both acoustic and seismic
amplitudes were largest during the opening sequence, declined rapidly after the
principal explosion and capsule-related signals, and increased again during the
dense Phase~III sequence. Later events were generally weaker than the
opening explosions, although several isolated impulses remained measurable late
in the record.

The principal explosion was the largest event by both pressure and vector PGV.
Several adjacent catalogue windows around the principal explosion also
contained part of its large waveform. They were retained in the chronological
catalogue but merged into one nonoverlapping event complex for energy
integration, preventing the same blast energy from being counted repeatedly.

Under the hemispherical-radiation model, the named principal-explosion window
contained approximately $1.89\times10^9$~J of acoustic energy, equivalent to
approximately 453~kg TNT as radiated acoustic energy
(Table~\ref{tab:named_energetics}). At an illustrative 1\% acoustic efficiency,
this would correspond to approximately 45~t TNT total yield; efficiencies of
0.1--10\% span approximately 4.5--450~t. These values are model scenarios, not
calibrated blast yields. The five named-window values in
Table~\ref{tab:named_energetics} are not summed directly with nearby catalogue
windows; cumulative energy was calculated only after all overlaps had been
merged into disjoint event complexes.

\begin{table*}[t]
\caption{Acoustic energetics of the five named opening-sequence windows.
Peak SPL is referenced to 20~$\mu$Pa. $E_{2\pi}$ assumes hemispherical
radiation without atmospheric attenuation, and acoustic TNT equivalent is
$E_{2\pi}/(4.184\times10^6~\mathrm{J\,kg^{-1}})$. The final column is an
illustrative total-yield scenario at 1\% acoustic efficiency, not a calibrated
blast yield. Nearby overlapping catalogue windows were merged before any
cumulative calculation, so these rows must not be added to overlapping raw
event-window estimates. Intended for the Supplementary Material.}
\label{tab:named_energetics}
\centering
\small
\begin{tabular}{lrrrrr}
\hline
Signal & Peak SPL (dB) & SEL (dB) & $E_{2\pi}$ (J) & Acoustic TNT (kg) & Yield at 1\% (t TNT) \\
\hline
Upper-stage event & 127.95 & 109.90 & $1.19\times10^6$ & 0.285 & 0.0285 \\
Principal explosion & 156.47 & 141.89 & $1.89\times10^9$ & 453 & 45.3 \\
Capsule impact & 142.45 & 126.26 & $5.18\times10^7$ & 12.4 & 1.24 \\
Capsule explosion & 143.30 & 126.97 & $6.10\times10^7$ & 14.6 & 1.46 \\
Event 024 & 142.06 & 125.39 & $4.24\times10^7$ & 10.1 & 1.01 \\
\hline
\end{tabular}
\end{table*}

The median dimensional pressure-to-PGV ratio for the full catalogue was
$4.58\times10^{5}\ \si{\pascal\per\metre\second}$. The ratio showed much more
scatter among events with vector PGV below approximately
$10^{-4}\ \si{\metre\per\second}$. This behavior coincided with decreasing
seismic SNR and reduced capture of the event maximum, indicating that at least
part of the increased low-amplitude scatter resulted from the seismic noise
floor rather than a demonstrable change in acoustic--seismic coupling.


\subsection{Acoustic pressure versus ground velocity}
\label{sec:results_regression}

Forty-six events met the joint acoustic-SNR, seismic-SNR, and window-capture
criteria for regression. Within this subset, peak-to-peak pressure and
vector PGV were strongly correlated in logarithmic space. The fitted power-law
exponent was
\begin{equation}
  b = 0.962,
\end{equation}
with a 95\% event-bootstrap interval of 0.893--1.009. The regression explained
95.3\% of the variance in $\log_{10}$ pressure ($R^2=0.953$), and the
Spearman correlation was 0.931.

The confidence interval included unit slope, supporting approximately
proportional pressure and ground-velocity amplitudes. Equal-weight orthogonal
regression gave a slope of 0.986; excluding the principal explosion gave 0.958
with $R^2=0.937$; and leave-one-out slopes ranged only from 0.952 to 0.971.
The near-unit relationship was therefore not controlled by one high-leverage
event. The median pressure-to-vector-PGV ratio in the selected subset was
$4.34\times10^{5}\ \si{\pascal\per\metre\second}$.

The largest events extended the general acoustic--seismic trend rather than
forming a separate population. Weak events showed greater scatter as seismic
SNR decreased and were excluded from physical interpretations of changing
energy partition.

\begin{figure*}[t]
  \centering
  \includegraphics[width=0.82\textwidth]{publication_pressure_pgv_regression.pdf}
  \caption{Peak-to-peak pressure versus three-component vector PGV for the 46
  events satisfying independent pressure-SNR, seismic-SNR, and window-capture
  criteria. The solid line is the log--log regression, shading is its bootstrap
  95\% interval, and the dashed line shows unit slope. The fitted slope is
  0.962 (95\% interval 0.893--1.009; $R^2=0.953$), consistent with
  approximately proportional pressure and ground-velocity amplitudes.}
  \label{fig:pressure_pgv}
\end{figure*}

\subsection{Frequency-dependent pressure--ground coupling}
\label{sec:results_transfer}

Pressure and vertical-velocity spectra were broadband but differed
systematically. Vertical response per unit pressure was weak below
approximately 3~Hz and enhanced in broad bands near 4--5 and 20--24~Hz, with a
weaker feature near 13--14~Hz (Figure~\ref{fig:transfer}). The event-based H1
magnitude reached approximately $5$--$7\times10^{-6}$
(m~s$^{-1}$)~Pa$^{-1}$ in the principal bands. Event coherence was
approximately 0.61--0.68 near 4.5--5.3~Hz and 0.43--0.70 near 21--28~Hz,
exceeding the frequency-specific permutation thresholds.

The continuous Phase~I estimate reproduced the main structure, with a maximum
near 4.7~Hz of approximately $1.1\times10^{-5}$
(m~s$^{-1}$)~Pa$^{-1}$ and values of approximately
$5$--$7\times10^{-6}$ (m~s$^{-1}$)~Pa$^{-1}$ near 21--25~Hz. The 13--14~Hz
feature was present but less consistent. Removing the strongest events,
equalizing pressure RMS, and changing pressure-beam construction preserved the
two principal bands. They are therefore interpreted as broad, repeatable
coupling bands rather than sharply resolved resonances. The empirical transfer
estimate describes the BCHH sensor--site configuration and is not a universal
air-to-ground transfer function.

\begin{figure*}[t]
  \centering
  \includegraphics[width=0.94\textwidth]{pressure_dhz_transfer_sensitivity.pdf}
  \caption{Empirical pressure-to-vertical-ground-velocity transfer estimates
  and sensitivity tests for the high-quality event ensemble and continuous
  Phase~I record. Bootstrap intervals describe sampling variability and
  permutation thresholds identify frequency-dependent coherence above
  mismatched pressure--velocity pairs. Enhanced response near 4--5 and
  20--24~Hz persists after excluding dominant events, equalizing pressure RMS,
  and changing pressure-beam construction. Event and continuous estimates
  overlap in time and are complementary rather than independent.}
  \label{fig:transfer}
\end{figure*}


\section{Discussion}
\label{sec:discussion}

\subsection{Independent forensic reconstruction}
\label{sec:discussion_reconstruction}

The BCHH observations provide an independent forensic record of the Falcon~9
accident. The original analysis was organized around five practical questions:
whether a detectable energetic precursor occurred; how many discrete events
followed and when; how large the principal events were; whether the impulsive
signals originated consistently from SLC-40; and whether their propagation
speeds were compatible with atmospheric airwaves. The seismic and infrasound
records address these questions without reliance on proprietary launch-vehicle
telemetry.

No anomalous impulsive precursor attributable to the accident was identified
during the preceding 37~h. The sequence subsequently produced 153 impulsive
infrasound events over approximately 26~min; a simple pulse morphology was
measurable for 152, while the principal explosion was retained as a complex
waveform. High-quality array solutions consistently pointed toward SLC-40 and
gave an apparent velocity of 351.98~m~s$^{-1}$, essentially identical to the
350.98~m~s$^{-1}$ meteorological prediction.

The combined seismic, infrasound, and publicly available video observations
permit a detailed reconstruction of the temporal evolution of the Falcon~9
explosion sequence (Figures~\ref{fig:opening} and~\ref{fig:multimodal}). Although the datasets record
different physical manifestations of the accident, they are mutually
consistent and together provide a more complete chronology than any individual
dataset alone.

The first visible flash was mapped to the upper-stage source time at
13:07:11.913~UTC. Organized low-frequency ground motion began approximately
1.7--2.0~s later, before the predicted ground-coupled airwave at +4.046~s, but
its slow apparent velocity and radial-plane polarization do not support an
unambiguous direct P-wave interpretation. No anomalous seismic or infrasonic
activity was detected during the preceding 37~h, indicating that the observed
sequence began abruptly.

Approximately 3.60~s after the upper-stage event, the principal explosion
produced by far the strongest acoustic and seismic signals recorded by BCHH.
Public video shows rapid expansion of the principal fireball and catastrophic
destruction of the launch vehicle. Several seconds later, acoustic and ground-coupled signals occurred during
the interval in which video shows the payload or capsule descending and
apparently reaching the ground. The current raw waveforms do not isolate an
unambiguous direct seismic impact phase, so the association remains tentative
and is based on the combined timing and video evidence rather than on a uniquely
identified impact phase.

Following the initial failure phase, activity evolved through three
additional phases of progressively smaller impulsive events separated by
intervals of greatly reduced activity. Between the first and second phases, a
prolonged continuous signal persisted for approximately 3.5~min. The signal is
prominent seismically and may also be present on DD2, but it does not appear as
a sequence of discrete impulsive infrasound events. A possible association with
sustained launch-pad activity is plausible, but its physical source is not
resolved by the present single-station observations.

Taken together, the observations demonstrate that the Falcon~9 accident was not a single explosion but a prolonged sequence of discrete energetic events extending over more than 26~min. The event catalogue provides a quantitative framework for examining that evolution and comparing it with subsequent engineering investigations.

\subsection{Comparison with the official failure investigation}
\label{sec:discussion_investigation}

At the time the original analyses were performed in 2016--2017, the engineering cause of the accident had not yet been established. SpaceX subsequently released results of its failure investigation, providing an independent chronology against which the geophysical observations can be evaluated.

According to that investigation, the initiating failure involved rapid release
of helium from a composite overwrapped pressure vessel immersed within the
second-stage liquid-oxygen tank. The resulting overpressure led to structural
failure of the second stage, followed by rapid destruction of the launch
vehicle and ignition of the propellant inventory. Although the geophysical
observations alone cannot identify these engineering mechanisms, they
independently document the timing and relative magnitudes of the principal
energetic events. The absence of a detectable preceding energetic signal is
consistent with an initiating process internal to the vehicle, but the BCHH
observations alone cannot identify that process or exclude every possible
external mechanism.

The initial infrasound event represents the first identifiable atmospheric
manifestation of the accident and is consistent in timing with the earliest
visible evidence of vehicle failure. Approximately 3.60~s later, the principal
explosion produced the largest acoustic and seismic signals of the sequence and
coincided with rapid growth of the fireball in public video. The close
correspondence between the independently derived geophysical chronology and
subsequently published engineering timeline supports the interpretation that
the event catalogue accurately captures the evolution of the accident.

Importantly, the event catalogue, propagation measurements, amplitude estimates, and timing relationships were derived from the seismic and infrasound recordings independently of the later engineering interpretation. The SpaceX findings therefore provide an external comparison rather than the basis for the geophysical measurements.

\subsection{Seismo-acoustic characteristics of a near-field rocket explosion}
\label{sec:discussion_characteristics}

The Falcon~9 dataset provides a rare opportunity to examine a catastrophic technological explosion at a source distance of only approximately 1.4~km. At this range, direct acoustic arrivals dominate, minimizing uncertainties associated with long-range atmospheric propagation that complicate regional infrasound observations.

The principal explosion generated positive pressures exceeding 1.3~kPa at
BCHH while producing much smaller ground-velocity amplitudes. Across the 46
events with reliable pressure and seismic measurements, however, pressure and
three-component PGV followed a nearly proportional relation over their observed
dynamic range (Figure~\ref{fig:pressure_pgv}). The near-unit exponent and its
stability to removal of the principal explosion indicate a repeatable local
air-to-ground response rather than two distinct populations of large and small
events. Smaller events were more readily detected in pressure than in ground
velocity because their seismic measurements approached the BCHH noise floor.

The frequency-domain results show that this response was not spectrally
uniform. The event ensemble and continuous Phase~I analysis both identified
broad coherent response bands near 4--5 and 20--24~Hz, with a weaker feature
near 13--14~Hz (Figure~\ref{fig:transfer}). These empirical transfer estimates
characterize the combined local ground, installation, and sensor response at
BCHH; they should not be interpreted as a universal acoustic-to-seismic
transfer function or as a direct measure of source energy partition.

Array measurements further demonstrate that the impulsive events were
remarkably consistent in both back azimuth and apparent propagation speed
despite their broad amplitude range (Figure~\ref{fig:array_summary}). The
approximately 1$^\circ$ systematic directional offset is smaller than can be
resolved reliably from cable-end sensor positions reconstructed to within only
1--2~m, and therefore does not establish source motion within the launch
complex.

Classical symmetric N-wave morphology was not observed, even for the
principal explosion. The available data do not uniquely distinguish among
source complexity, near-ground propagation, waveform evolution over the
1.4-km path, and instrumental or local-site effects, so no single explanation
is preferred.

\subsection{Planar and finite-range propagation}
\label{sec:discussion_propagation}

The modern Python analysis extends the original planar-wave treatment by explicitly evaluating a finite-range circular wavefront. At a source distance of approximately 1.4~km and an array aperture of approximately 30~m, wavefront curvature is small but potentially measurable. Comparison of the planar and fixed-source circular models provides a direct test of whether curvature affects the inferred propagation speed or coherence.

The fixed-SLC-40 model preferred a median speed 4.0~m~s$^{-1}$ greater than the
planar solution while reducing coherence by a median of 0.00512. This small
degradation shows that explicitly imposing the nominal pad position did not
improve waveform alignment. The speed increase is best understood as
compensation for the approximately 1$^\circ$ difference between the measured
wavefront direction and the direction imposed by the nominal coordinates, not
as evidence for faster atmospheric propagation.

The finite-range sensitivity calculations also clarify the resolution of the
three-element array (Figure~\ref{fig:finite_distance}). Two independent
relative delays are sufficient to constrain direction and apparent speed under
a planar model, but they cannot independently determine source easting, source
northing, and propagation speed. Every along-range profile remained open at
the search boundary for a coherence decrease of 0.01. The open along-range
ridges therefore represent a fundamental range-resolution limitation rather
than a failure of the optimization procedure, and the simpler planar result is
preferred.

\subsection{Comparison with other technological and natural explosions}
\label{sec:discussion_comparison}

Large technological explosions have previously been investigated using
seismic and infrasound observations, including the 2005 Buncefield oil-depot
explosion \citep{ottemoller_2008,green_2008,ceranna_2009}, the 2014 Moerdijk
chemical-plant explosions \citep{smink_2019}, and the 2020 Beirut
ammonium-nitrate explosion
\citep{pilger_beirut_2021,kim_pasyanos_2022}. Such studies demonstrate the value of seismo-acoustic observations for estimating source location, explosive yield, and source characteristics. Most, however, rely primarily on regional observations, where atmospheric propagation strongly influences arrival times, amplitudes, and waveforms and detailed propagation modelling is required.

The Falcon~9 dataset differs in several respects. BCHH was located only approximately 1.4~km from the launch pad, so direct acoustic and seismic arrivals dominate and propagation uncertainties are comparatively small. Rather than emphasizing yield estimation alone, the proximity of the observations permits reconstruction of a complex sequence comprising more than 150 individual impulsive events. Public video and the subsequent failure investigation provide independent constraints against which the geophysical chronology can be evaluated. Together, these datasets make the accident an unusually well-constrained benchmark for near-field seismo-acoustic analysis of a technological explosion.

In several respects, the prolonged sequence resembles sustained explosive activity during volcanic eruptions, in which a major explosion is followed by numerous smaller impulses over minutes to tens of minutes. The event-catalogue approach therefore has applications beyond rocket launches, providing a common framework for analysing complex atmospheric explosion sequences and integrating seismic, infrasound, and optical observations to reconstruct source evolution.

\subsection{Implications for rocket-launch monitoring}
\label{sec:discussion_monitoring}

The accidental recording of the Falcon~9 failure demonstrates the forensic
value of continuous seismo-acoustic monitoring at launch facilities. A single
well-calibrated station constrained the onset, absence of detectable energetic
precursors, number and timing of subsequent impulsive events, relative
amplitudes, source direction, and propagation characteristics without access to
proprietary engineering telemetry.

Since 2016, launch cadence and the number of launch providers at Kennedy Space Center and Cape Canaveral have increased substantially. The methods developed here---including systematic event cataloguing, array-based propagation analysis, and integration with independent observational datasets---provide a foundation for monitoring both nominal launches and off-nominal events.

Future networks employing multiple seismo-acoustic arrays could substantially improve source-location capability while enabling more detailed characterization of explosion dynamics, sonic booms, and other transient atmospheric phenomena associated with rocket operations.

\subsection{Limitations}
\label{sec:discussion_limitations}

Several limitations should be considered. The observations were obtained from a single three-element seismo-acoustic array, preventing independent triangulation of individual sources. The array analysis therefore provides robust directional and apparent-speed constraints but limited source-range resolution.

Event identification was performed manually because of the modest data volume and complexity of the sequence. Although this produced an internally consistent catalogue, larger datasets would benefit from automated detection algorithms trained and validated against the manual reference catalogue.

The infrasound calibrations were empirical field estimates rather than
laboratory-traceable absolute calibrations, and their absolute uncertainty
requires final documentation. Public-video timing is also limited by uncertain
camera locations, frame timing, and absent absolute timestamps in some
recordings.

The forensic conclusions are constrained by the types of processes detectable
seismo-acoustically. The absence of a precursor means that no anomalous
energetic seismic or infrasonic signal was identified above the effective noise
and detection limits of BCHH. It does not exclude slow structural degradation,
internal fluid or mechanical processes, weak sources, or other mechanisms that
would not generate a detectable signal at the array.

Interpretations of payload or capsule impact and of the prolonged continuous
signal remain provisional. The capsule-event labels are supported by relative
timing and video context, but no unique direct seismic impact phase was
identified. The continuous signal is well resolved seismically, yet a single
station cannot determine whether it represents sustained combustion, shallow
ground motion, or another launch-complex process.

Finally, the acoustic and seismic energy estimates should be regarded as
first-order approximations because of uncertainties in calibration,
atmospheric attenuation, near-surface properties, source directivity, and the
assumption of hemispherical radiation. Overlapping catalogue windows were
merged before cumulative acoustic energy was calculated, but the cumulative
value includes only the union of high-quality and named-event intervals.
TNT-equivalent total yields additionally depend on an assumed acoustic
efficiency and are reported as sensitivity scenarios, not calibrated blast
yields. These limitations do not preclude relative comparisons within the
sequence but must be considered when comparing absolute energies with other
events.

\section{Conclusions}
\label{sec:conclusions}

The 1 September 2016 Falcon~9 static-fire accident produced one of the most comprehensively documented near-field technological explosions recorded by a colocated seismic and infrasound array. Continuous observations acquired at BCHH, approximately 1.4~km from SLC-40, captured the complete 26-minute sequence with high signal-to-noise ratio, providing a rare opportunity to investigate the temporal evolution of a catastrophic launch-vehicle failure using geophysical observations.

Manual association identified a catalogue of 153 impulsive infrasound events.
Of these, 152 yielded valid positive-compression/negative-rarefaction
measurements; the principal explosion was retained as a complex noncanonical
event. The records also contain a prolonged interval of continuous seismic
tremor that is not counted as an additional discrete event. The catalogue is
the principal scientific product of this study and provides quantitative
measurements of event timing, amplitude, array coherence, apparent propagation
speed, and back azimuth (Figures~\ref{fig:overview} and
\ref{fig:array_summary}).

The observations address the principal forensic questions that motivated the
original analysis. No anomalous energetic seismic or infrasonic precursor
attributable to the accident was identified during the preceding 37~h. This
result constrains detectable impulsive precursors but does not exclude
non-energetic or seismo-acoustically undetectable initiating processes. The
accident generated 153 impulsive infrasound events over approximately 26~min,
rather than a single explosion. Approximately 3.60~s after the initial
upper-stage signal, the principal explosion produced the largest pressure and
ground-motion amplitudes of the sequence (Table~\ref{tab:principal_events}).
High-quality array solutions consistently pointed toward SLC-40. Their
coherence-weighted apparent speed was 351.98~m~s$^{-1}$, agreeing with the
independently predicted meteorological effective speed of
350.98~m~s$^{-1}$.

Pressure and three-component PGV were nearly proportional among 46 events with
reliable measurements, with a log--log slope of 0.962 and $R^2=0.953$
(Figure~\ref{fig:pressure_pgv}). Event and continuous analyses within the same
Phase~I record identified broad coherent air-to-ground response bands near
4--5 and 20--24~Hz (Figure~\ref{fig:transfer}). In addition, organized
low-frequency ground motion began 1.7--2.0~s after the upper-stage source time
and before the predicted airwave. Its 0.71--0.84~km~s$^{-1}$ apparent speed and
radial-plane polarization favor a slow shallow-ground or surface-wave packet
over an unambiguous direct P wave (Figure~\ref{fig:preairwave}).

Comparison with public video and the subsequently published SpaceX failure investigation shows that the independently derived geophysical chronology is consistent with the documented evolution of the accident. The observations therefore provide an independent forensic record without relying on proprietary engineering telemetry.

More broadly, the study demonstrates the value of continuous near-field
seismo-acoustic monitoring at major launch facilities. Even a single
well-calibrated array was sufficient to constrain the timing, evolution,
relative magnitudes, source direction, and propagation characteristics of a
complex launch-vehicle failure. Multiple arrays would add source-location
capability and improve characterization of nominal launches, anomalous events,
sonic booms, and explosion dynamics.

By making the waveform archive, event catalogue, and analysis software openly available, this study will establish a benchmark dataset for developing and evaluating future seismo-acoustic methods and for comparison with subsequent rocket-launch accidents and other complex atmospheric explosions.

\begin{acknowledgements}
% TODO: confirm acknowledgements, funding, KSC support, array deployment assistance, weather-data providers, SpaceX wording, and Astronaut Beach House access.
Acknowledgements to be completed.
\end{acknowledgements}

\section*{Data and code availability}

% TODO: replace placeholders with persistent repository citations and DOIs.
Continuous waveform data, station metadata, manual arrival picks, the event catalogue, Python notebooks, archived MATLAB routines, derived tables, and figure-generation scripts will be deposited in an open repository with a persistent identifier. The repository will identify which workflow generated each reported product. GitHub-hosted development code will be archived as a versioned release in a persistent repository such as Zenodo.

\section*{Competing interests}

The authors declare that they have no competing interests.

% Requires mybibfile.bib in the manuscript directory.
\bibliography{mybibfile}

\end{document}
